% 数学ノート集・共通テンプレート v2（積分論2026年度第1部の版面を参考）
% A5判。章内で定義・定理・例・注意等の番号を共通化。
% コンパイル： platex -interaction=nonstopmode math_note_template_v2.tex (2回)
%              dvipdfmx math_note_template_v2.dvi
\documentclass[a5paper,10pt,onecolumn]{jsarticle}
\usepackage[dvipdfmx,top=22mm,bottom=20mm,inner=19mm,outer=18mm,headheight=13pt,headsep=6mm]{geometry}
\usepackage{amsmath,amssymb,amsthm,mathtools}
\usepackage{enumitem,fancyhdr,titlesec}
\usepackage[dvipdfmx,bookmarks=true,colorlinks=false]{hyperref}
\usepackage{pxjahyper}

% ----- 差し替える文書情報 -----
\newcommand{\NoteTitle}{微分形式とPoincar\'eの補題}
\newcommand{\NoteSeries}{数学ノート集}

% ----- 文字・組版 -----
% 参考PDFの質素な明朝＋Computer Modernの印象を引き継ぎながら，
% A5判で窮屈にならないよう本文の行間を少し確保する．
\linespread{1.13}
\setlength{\parindent}{1zw}
\setlength{\parskip}{0pt}
\setlength{\abovedisplayskip}{6pt plus 2pt minus 2pt}
\setlength{\belowdisplayskip}{6pt plus 2pt minus 2pt}
\setlength{\abovedisplayshortskip}{3pt plus 2pt minus 1pt}
\setlength{\belowdisplayshortskip}{4pt plus 2pt minus 1pt}
\setlist[enumerate]{topsep=3pt,itemsep=1pt,parsep=0pt,leftmargin=2.4zw}
\titleformat{\section}{\normalfont\normalsize\bfseries}{\thesection}{0.8zw}{}
\titlespacing*{\section}{0pt}{2.0ex plus .4ex}{.9ex}
\titleformat{\subsection}{\normalfont\normalsize\bfseries}{\thesubsection}{0.7zw}{}
\titlespacing*{\subsection}{0pt}{1.5ex plus .3ex}{.6ex}
\raggedbottom

% ----- 柱とページ番号：ページ上部、罫線なし -----
\fancypagestyle{mathnote}{%
  \fancyhf{}%
  \fancyhead[L]{\footnotesize \NoteSeries\ （\NoteTitle）}%
  \fancyhead[R]{\footnotesize \thepage}%
  \renewcommand{\headrulewidth}{0pt}%
  \renewcommand{\footrulewidth}{0pt}%
}
\pagestyle{mathnote}

% ----- 節内通し番号の定義・定理群 -----
\newtheoremstyle{mathnotetheorem}
 {4.5pt}{4.5pt}{\normalfont}{}%
 {\bfseries}{}{0.65zw}%
 {\thmname{#1}\thmnumber{ #2}\thmnote{（#3）}}
\theoremstyle{mathnotetheorem}
\newtheorem{theorem}{定理}[section]
\newtheorem{definition}[theorem]{定義}
\newtheorem{lemma}[theorem]{補題}
\newtheorem{proposition}[theorem]{命題}
\newtheorem{corollary}[theorem]{系}
\newtheorem{example}[theorem]{例}
\newtheorem{remark}[theorem]{注意}
\renewcommand{\proofname}{証明}
\renewcommand{\qedsymbol}{$\square$}
\makeatletter
\renewenvironment{proof}[1][\proofname]{%
  \par\pushQED{\qed}\normalfont\topsep4\p@\@plus2\p@\relax
  \trivlist\item[\hskip\labelsep\bfseries #1]\ignorespaces
}{\popQED\endtrivlist\@endpefalse}
\makeatother

% ----- 使うノートで適宜拡張する記号 -----
\newcommand{\R}{\mathbb{R}}
\newcommand{\norm}[1]{\lVert #1\rVert}
\newcommand{\abs}[1]{\lvert #1\rvert}

\begin{document}
\section*{読了後にできること}
本稿を読むための前提は，微分形式，外微分，Poincaréの補題．
本稿の具体的な到達課題は，閉形式の原始関数を具体的に構成する．

\section{交代多重線形形式}

本稿では $U\subset\mathbb R^n$ を開集合とし，
微分形式は特に断らない限り $C^\infty$ 級とする．

\begin{definition}
有限次元実線形空間 $V$ 上の交代 $k$-重線形形式全体を
$\Lambda^kV^*$ と書く．$k=0$ のとき
$\Lambda^0V^*=\mathbb R$ と約束する．
交代形式 $\alpha\in\Lambda^kV^*$ と
$\beta\in\Lambda^\ell V^*$ の\emph{外積}を
\[
 \begin{aligned}
 &(\alpha\wedge\beta)(v_1,\dots,v_{k+\ell})\\
 &=\frac1{k!\ell!}\sum_{\sigma\in S_{k+\ell}}
   \operatorname{sgn}(\sigma)\,
   \alpha(v_{\sigma(1)},\dots,v_{\sigma(k)})\\
 &\qquad\cdot
   \beta(v_{\sigma(k+1)},\dots,v_{\sigma(k+\ell)}).
 \end{aligned}
\]
によって定める．
\end{definition}

\begin{lemma}[外積の基本性質]\label{lem:wedge}
外積は結合的であり，
\[
 \alpha\wedge\beta=(-1)^{k\ell}\beta\wedge\alpha
\]
を満たす．
また $V=\mathbb R^n$ の双対基底
$dx^1,\dots,dx^n$ に対して
$dx^{i_1}\wedge\cdots\wedge dx^{i_k}$
（$1\le i_1<\cdots<i_k\le n$）は
$\Lambda^kV^*$ の基底をなす．
\end{lemma}
\begin{proof}
外積の定義は交代化したテンソル積である．
$\alpha,\beta,\gamma$ の外積を三つまとめて展開すると，
いずれの結合の仕方でも
全ての $(k,\ell,m)$-shuffleについて
同じ符号付き和になるため結合的である．
$\alpha,\beta$ のブロックを入れ替える置換には
$k\ell$ 回の隣接交換が必要なので
符号は $(-1)^{k\ell}$ となり，
次数付き交換法則を得る．
最後に任意の交代 $k$-形式は基底ベクトル
$(e_{i_1},\dots,e_{i_k})$ 上の値によって一意に決まり，
重複添字を持つ値は零で，
添字の交換による符号が決まる．
従って上の外積全体が基底をなす．
\end{proof}

\section{微分形式と外微分}

\begin{definition}
$U$ 上の\emph{$k$-次微分形式}とは，
各 $x\in U$ に交代 $k$-形式
$\omega_x\in\Lambda^k(\mathbb R^n)^*$ を滑らかに対応させるものをいう．
その全体を $\Omega^k(U)$ と書く．
座標では
\[
 \omega=\sum_{i_1<\cdots<i_k}
 a_{i_1\cdots i_k}(x)\,
 dx^{i_1}\wedge\cdots\wedge dx^{i_k}
\]
と一意に書ける．
\end{definition}

\begin{definition}
$\omega=\sum_I a_I\,dx^I\in\Omega^k(U)$ に対し，
\[
 d\omega=\sum_I\sum_{j=1}^n
 \frac{\partial a_I}{\partial x^j}
 dx^j\wedge dx^I
\]
と定める．これを\emph{外微分}という．
\end{definition}

\begin{proposition}[外微分の基本性質]\label{prop:exterior}
$\omega\in\Omega^k(U)$，
$\eta\in\Omega^\ell(U)$ とすると
\[
 d(\omega\wedge\eta)
 =d\omega\wedge\eta+(-1)^k\omega\wedge d\eta,
 \qquad d(d\omega)=0.
\]
また $d$ は線形である．
\end{proposition}
\begin{proof}
線形性は定義から直ちに従う．
$\omega=a\,dx^I$，$\eta=b\,dx^J$ の場合，
\[
 d(\omega\wedge\eta)
 =d(ab)\wedge dx^I\wedge dx^J
 =(b\,da+a\,db)\wedge dx^I\wedge dx^J.
\]
一方
$d\omega=da\wedge dx^I$，
$d\eta=db\wedge dx^J$ である．
$dx^I$ の次数は $k$ なので
$dx^I\wedge db=(-1)^kdb\wedge dx^I$
となり，Leibniz則を得る．
一般の場合は線形性で従う．

また
\[
 d^2(a\,dx^I)
 =\sum_{i,j}\partial_j\partial_i a\,
 dx^j\wedge dx^i\wedge dx^I.
\]
$i,j$ を交換すると偏導関数の値は等しく，
外積の符号は逆になる．$i=j$ の項は零なので，
全ての項が打ち消し合い $d^2=0$ となる．
\end{proof}

\begin{definition}
$d\omega=0$ を満たす形式を\emph{閉形式}，
$\omega=d\eta$ と書ける形式を\emph{完全形式}という．
$d^2=0$ によって任意の完全形式は閉形式である．
\end{definition}

\section{写像による引き戻し}

\begin{definition}
$U\subset\mathbb R^n$，
$V\subset\mathbb R^m$ を開集合，
$F:U\to V$ を滑らかな写像とする．
$\omega\in\Omega^k(V)$ に対し
\[
 (F^*\omega)_x(v_1,\dots,v_k)
 =\omega_{F(x)}(DF_xv_1,\dots,DF_xv_k)
\]
と定め，これを $\omega$ の\emph{引き戻し}という．
\end{definition}

\begin{theorem}[引き戻しと外微分]\label{thm:pullback}
滑らかな写像 $F:U\to V$ と
$\omega,\eta$ に対し
\[
 F^*(\omega\wedge\eta)=F^*\omega\wedge F^*\eta,
 \qquad
 d(F^*\omega)=F^*(d\omega)
\]
が成り立つ．
\end{theorem}
\begin{proof}
外積については，その交代和の各項で
$DF_x$ を各変数に作用させればよく，
定義を展開して直接一致する．
外微分については，座標関数 $y^j$ を用いて
$\omega=\sum_I a_I(y)\,dy^{i_1}\wedge\cdots\wedge dy^{i_k}$
と書く．
\[
 F^*\omega=\sum_I(a_I\circ F)\,
 dF^{i_1}\wedge\cdots\wedge dF^{i_k}
\]
である．
$F^{i_j}$ は関数なので $d^2F^{i_j}=0$ であり，
Leibniz則から
\[
 d(F^*\omega)=
 \sum_I d(a_I\circ F)\wedge
 dF^{i_1}\wedge\cdots\wedge dF^{i_k}.
\]
実多変数の連鎖律によって
$d(a_I\circ F)=\sum_j(\partial_ja_I\circ F)\,dF^j
=F^*(da_I)$ だから，
右辺は $F^*(d\omega)$ に等しい．
\end{proof}

\section{星状領域とホモトピー作用素}

\begin{definition}
開集合 $U\subset\mathbb R^n$ が\emph{星状}であるとは，
ある $a\in U$ が存在して
$a+t(x-a)\in U$ が全ての $x\in U$，
$0\le t\le1$ に対して成り立つことをいう．
\end{definition}

以下，座標を平行移動して星状中心を $0$ と仮定する．
\begin{definition}
$k\ge1$ と $\omega\in\Omega^k(U)$ に対して
\[
 (K\omega)_x(v_1,\dots,v_{k-1})
 =\int_0^1t^{k-1}
 \omega_{tx}(x,v_1,\dots,v_{k-1})\,dt
\]
と定める．$K:\Omega^k(U)\to\Omega^{k-1}(U)$ を
\emph{ホモトピー作用素}という．
$K:\Omega^0(U)\to0$ と約束する．
\end{definition}

\begin{lemma}\label{lem:k-smooth}
$K\omega$ は滑らかな $(k-1)$-形式である．
\end{lemma}
\begin{proof}
任意のコンパクト集合 $C\subset U$ を取ると，
$\{tx:x\in C,\ 0\le t\le1\}$ は $U$ のコンパクト部分集合となる．
$\omega$ の係数はその集合上で全ての階数について有界であり，
微分も連続である．
従って $t\in[0,1]$ に関する積分と
$x$ に関する任意階の微分を
コンパクト集合上で交換できる．
これにより $K\omega$ の係数関数は滑らかである．
交代性は $\omega$ の交代性から従う．
\end{proof}

\begin{theorem}[ホモトピー公式]\label{thm:homotopy-formula}
$U$ を $0$ を中心とする星状開集合とする．
任意の $\omega\in\Omega^k(U)$ に対して
\[
 d(K\omega)+K(d\omega)=\omega-c_0^*\omega
\]
が成り立つ．
ここで $c_0:U\to U$ は定数写像 $c_0(x)=0$ であり，
$k\ge1$ なら $c_0^*\omega=0$，
$k=0$ なら $c_0^*f=f(0)$ とする．
\end{theorem}
\begin{proof}
$H:[0,1]\times U\to U$ を
$H(t,x)=tx$ とする．
任意の $k$-形式 $\omega$ の引き戻しは
一意に
\[
 H^*\omega=dt\wedge\eta_t+\theta_t
\]
と分解できる．ここで $\eta_t$ は $U$ 上の
$(k-1)$-形式，$\theta_t$ は $k$-形式であり，
各係数は $t$ に滑らかに依存する．
$\eta_t$ の定義を展開すると
\[
 (\eta_t)_x(v_1,\dots,v_{k-1})
 =t^{k-1}\omega_{tx}(x,v_1,\dots,v_{k-1})
\]
となるため
$K\omega=\int_0^1\eta_t\,dt$ である．

$[0,1]\times U$ の外微分を $d_{t,x}$，
$U$ 方向の外微分を $d_x$ と書く．
定義とLeibniz則から
\[
 d_{t,x}(H^*\omega)
 =dt\wedge(\partial_t\theta_t-d_x\eta_t)
  +d_x\theta_t.
\]
一方，定理\ref{thm:pullback}により
$d_{t,x}(H^*\omega)=H^*(d\omega)$ である．
右辺の $dt$ を含む成分は，
$K(d\omega)$ の定義に現れるので，
\[
 K(d\omega)=\int_0^1(\partial_t\theta_t-d_x\eta_t)\,dt.
\]
従って
\begin{align*}
 d(K\omega)+K(d\omega)
 &=d_x\int_0^1\eta_t\,dt
   +\int_0^1(\partial_t\theta_t-d_x\eta_t)\,dt\\
 &=\theta_1-\theta_0\\
 &=H_1^*\omega-H_0^*\omega
 =\omega-c_0^*\omega.
\end{align*}
最後に $H_1=\operatorname{id}_U$，
$H_0=c_0$ を用いた．
\end{proof}

\section{Poincar\'eの補題}

\begin{theorem}[Poincar\'eの補題]\label{thm:poincare}
$U\subset\mathbb R^n$ を星状開集合とする．
$k\ge1$ と $\omega\in\Omega^k(U)$ について
$d\omega=0$ なら，
ある $\eta\in\Omega^{k-1}(U)$ が存在して
$\omega=d\eta$ となる．
すなわち，星状領域上の閉形式は完全である．
\end{theorem}
\begin{proof}
星状中心を平行移動して $0$ とする．
定理\ref{thm:homotopy-formula}を適用すると，
$k\ge1$ なので $c_0^*\omega=0$ であり，
また $d\omega=0$ だから
\[
 d(K\omega)=\omega-K(d\omega)=\omega
\]
となる．
従って $\eta=K\omega$ と置けば
$\omega=d\eta$ を満たす．
\end{proof}

\begin{corollary}[局所Poincar\'eの補題]\label{cor:local-poincare}
滑らかな多様体 $M$ 上の閉 $k$-形式
$\omega$（$k\ge1$）は局所的に完全である．
すなわち任意の $p\in M$ に対し，
$p$ の開近傍 $V$ と
$(k-1)$-形式 $\eta$ が存在して
$\omega|_V=d\eta$ となる．
\end{corollary}
\begin{proof}
$p$ の座標近傍 $(W,\varphi)$ を取る．
$\varphi(p)$ を中心とする十分小さい開球
$B\subset\varphi(W)$ を選び，$V=\varphi^{-1}(B)$ とする．
$B$ は星状なので
$\varphi^{-1}$ による $\omega$ の座標表示は
$B$ 上の閉形式であり，
定理\ref{thm:poincare}によって
ある形式の外微分となる．
その原始形式を $\varphi$ で引き戻せば
$V$ 上の $\eta$ を得る．
外微分と引き戻しの可換性は定理\ref{thm:pullback}による．
\end{proof}

\section{閉形式と完全形式の相違}

\begin{example}[穴を持つ領域上の閉形式]\label{ex:angular}
$U=\mathbb R^2\setminus\{0\}$ 上の1形式
\[
 \omega=\frac{-y\,dx+x\,dy}{x^2+y^2}
\]
は閉形式であるが，完全形式ではない．
\end{example}
\begin{proof}
$P(x,y)=-y/(x^2+y^2)$，
$Q(x,y)=x/(x^2+y^2)$ と置く．
直接計算すると
\[
 \partial_xQ
 =\frac{y^2-x^2}{(x^2+y^2)^2}
 =\partial_yP
\]
なので
$d\omega=(\partial_xQ-\partial_yP)\,dx\wedge dy=0$
となる．
一方，単位円を
$\gamma(t)=(\cos t,\sin t)$（$0\le t\le2\pi$）
とパラメータ表示すると
\[
 \gamma^*\omega
 =\frac{-\sin t(-\sin t)+\cos t\cos t}
       {\cos^2t+\sin^2t}\,dt=dt
\]
だから
$\int_\gamma\omega=2\pi$ である．
もし $\omega=df$ なら
曲線積分の基本定理によって
$\int_\gamma df=f(\gamma(2\pi))-f(\gamma(0))=0$
となるはずであり，矛盾する．
\end{proof}

\section{de Rhamコホモロジーへの導入}

\begin{definition}
$U\subset\mathbb R^n$ に対して
\[
 \begin{aligned}
 Z^k(U)&=\ker(d:\Omega^k(U)\to\Omega^{k+1}(U)),\\
 B^k(U)&=\operatorname{im}(d:\Omega^{k-1}(U)\to\Omega^k(U)).
 \end{aligned}
\]
と置く．$d^2=0$ によって
$B^k(U)\subset Z^k(U)$ なので，
\[
 H_{\mathrm{dR}}^k(U)=Z^k(U)/B^k(U)
\]
を定め，$k$ 次de Rhamコホモロジー群という．
\end{definition}

\begin{theorem}[星状領域のde Rhamコホモロジー]\label{thm:star-cohom}
$U\subset\mathbb R^n$ を空でない星状開集合とする．
このとき
\[
 H_{\mathrm{dR}}^0(U)\cong\mathbb R,\qquad
 H_{\mathrm{dR}}^k(U)=0\quad(k\ge1)
\]
が成り立つ．
\end{theorem}
\begin{proof}
$k\ge1$ ではPoincar\'eの補題によって
$Z^k(U)=B^k(U)$ となるので商は零である．
$k=0$ では $B^0(U)=0$ と約束する．
$Z^0(U)$ は $df=0$ を満たす滑らかな実数値関数全体である．
星状領域は道連結であり，
$df=0$ なら任意の道に沿う微分が零なので
$f$ は定数関数である．
従って $Z^0(U)$ は定数関数全体に等しく
$\mathbb R$ と同型である．
\end{proof}

\section{複素解析への接続}

\begin{proposition}[調和共役と閉1形式]\label{prop:harmonic-conjugate}
$U\subset\mathbb R^2$ を星状開集合とし，
$u\in C^\infty(U,\mathbb R)$ が
$\Delta u=u_{xx}+u_{yy}=0$ を満たすとする．
このとき $v\in C^\infty(U,\mathbb R)$ が存在して
$u+iv$ は正則となる．
$u+iv$ が正則となる $v$ は加法定数を除いて一意である．
\end{proposition}
\begin{proof}
\[
 \omega=-u_y\,dx+u_x\,dy
\]
と置く．外微分を計算すると
\[
 d\omega=(u_{xx}+u_{yy})\,dx\wedge dy=0
\]
である．定理\ref{thm:poincare}から
$\omega=dv$ となる滑らかな実数値関数 $v$ が存在する．
従って
$v_x=-u_y$，$v_y=u_x$ となり，
$u+iv$ はCauchy--Riemann方程式を満たす．
$u,v$ は滑らかだから
実微分と複素微分の同値性により $u+iv$ は正則である．
二つの調和共役の差は勾配零の関数となる．
$U$ は道連結なので，差は定数である．
\end{proof}

\begin{remark}
Poincar\'eの補題は閉形式が局所的には完全であることを主張するが，
大域的な完全性は領域の位相に依存する．
de Rhamコホモロジーはこの障害を代数的に記録する．
\end{remark}


\section{到達課題の完全解答}
\noindent\textbf{閉1形式の原始関数を計算する．}
\begin{proposition}
$\mathbb R^2$ 上の1形式
$\omega=y\,dx+x\,dy$ は閉形式であり，
$\omega=d(xy)$ なので完全形式でもある．
\end{proposition}
\begin{proof}
$d(y\,dx)=dy\wedge dx$，
$d(x\,dy)=dx\wedge dy$ より
\[
d\omega=dy\wedge dx+dx\wedge dy=0.
\]
一方積の微分公式から
$d(xy)=y\,dx+x\,dy=\omega$ である．
\end{proof}
\noindent\textbf{到達確認．}
$\eta=(2x+y)\,dx+x\,dy$ の原始関数を求めよ．
\par\noindent\textbf{解答．}
$d(x^2+xy)=(2x+y)\,dx+x\,dy$ だから
$x^2+xy+C$ が原始関数である．
\end{document}